\scanpage{112}

On a

\[c(f_V)\overset{\mathrm{dfn}}=w_2(f_V)+w_1(f_V)^2=(\unsure{\eta_0|y_0}+\eta_1+\eta_2)^2+(\eta_0\eta_1+\eta_1\eta_2+\eta_2\eta_0)\]

(Son annulation est néc et suff pour que le fibré en droites projectives envisagé soit isomorphe à celui défini \deleted{par \unsure{l'équat.|l'équation}} \inserted{par la} forme $x^2+y^2+z^2$). On a

\[\operatorname{Br}(f_V)\overset{\mathrm{dfn}}=c(f_V)-\underbrace{c(\deleted{\unsure{f_{2E}}}\,X^2+yz)}_{\varepsilon^2}=(\varepsilon+\eta_0+\eta_1+\eta_2)^2+(\eta_0\eta_1+\eta_1\eta_2+\eta_2\eta_0)\]

L'annulation de cette expression est néc et suff pour que le fibré en droites projectives \deleted{\unsure{envis.|associé}} défini par $f_V$ soit associable à un fibré vect. de rang 2, en fait, c'est la classe de Brauer de ce fibré. -- Le calcul universel se fait sur

\[S_0=\operatorname{Spec}\mathbb Z[\beta_0,\beta_1]_{2\beta_0\beta_1(4-\beta_0-\beta_1)}\]

On \unsure{notera|note} que dans $H^1(S_0,\mathbb Z/2)$, les \deleted{sections} \inserted{éléments} $\varepsilon,\eta_0,\eta_1,\eta_2$ forment une base. Il est plausible que les 10 expressions $\varepsilon^2$, $\eta_0^2$, $\eta_1^2$, $\eta_2^2$, $\varepsilon \eta_0$, $\varepsilon \eta_1$, $\varepsilon \eta_2$, $\eta_0\eta_1$, $\eta_1\eta_2$, $\eta_2\eta_0$ sont

% archive sheet number: 111
